\documentclass{article}
\usepackage{amsmath,amssymb,amsfonts}
\usepackage{physics}
\usepackage{hyperref}
\usepackage{tikz-cd}
\usepackage{geometry}
\geometry{margin=1in}

\title{Symmetry and Conservation Laws in Modern Physics}
\author{整理自知乎用户 X-yin 的回答}
\date{}

\begin{document}
\maketitle

\section{对称性与守恒律}
\begin{itemize}
  \item \textbf{时间平移不变性} $\Rightarrow$ 能量守恒：自然规律不依赖具体时刻。
  \item \textbf{空间平移不变性} $\Rightarrow$ 动量守恒：物理定律在空间中各处一致。
  \item \textbf{空间旋转不变性} $\Rightarrow$ 角动量守恒：没有绝对方向，旋转对称。
\end{itemize}

\section{从对称性到几何结构}
\begin{itemize}
  \item Lorentz 变换是 Minkowski 空间中的伪旋转。
  \item 广义相对论是局域化的平移不变理论，对应引力的几何化。
  \item Maxwell 理论是局域化的 $U(1)$ 不变性理论，对应电磁相互作用。
  \item Yang-Mills 理论是局域内部对称性（如 $SU(2)$, $SU(3)$）的推广。
  \item 超对称是超空间中沿费米方向的平移变换。
\end{itemize}

\section{规范场与力的来源}
\begin{itemize}
  \item 所有自然力来自\textbf{连续、局域的对称性}。
  \item 分立对称性（如时间反演、宇称）不产生力。
  \item 局域化超对称性 $\Rightarrow$ 超引力理论，包含引力。
\end{itemize}

\section{标准模型核心公式}
\[ Q = T_3 + \frac{Y}{2} \]
其中 $Q$ 是电荷，$T_3$ 是弱同位旋第三分量，$Y$ 是超荷。

\section{几何语言中的曲率与挠率}
\begin{align*}
  \mathfrak{T} &= \mathscr{D} e \quad \text{(挠率)} \\
  \mathfrak{R} &= \mathscr{D}^2 e \quad \text{(曲率)}
\end{align*}

\section{上同调结构：$Q^2 = 0$}
\subsection*{广义相对论}
\[ Q = \partial, \quad Q^2 = 0 \Rightarrow \text{转矩守恒} \]
\[ \star \int_{\partial\partial \Omega}(\mathscr{P} \wedge \mathfrak{R}) \equiv 0 \]

\subsection*{电动力学}
\[ Q = d \text{ 或 } \delta \Rightarrow \delta \boldsymbol{J} = 0 \]

\[ \boldsymbol{F} = d\boldsymbol{A}, \quad d\boldsymbol{F} = 0, \quad d\star \boldsymbol{F} = 4\pi \star \boldsymbol{J}, \quad \delta = \star d\star \]

\begin{tikzcd}
  \boldsymbol{A} \arrow[d] \\
  \boldsymbol{F} = d\boldsymbol{A} \arrow[r, ">>", "\star"] \arrow[d] & \star\boldsymbol{F} \arrow[d] \\
  d\boldsymbol{F} = 0 & d\star\boldsymbol{F} = 4\pi \star\boldsymbol{J} \arrow[r, ">> \delta = \star d \star"] & \delta \boldsymbol{F} = \boldsymbol{J} \\
  d^2 = 0 & d\star \boldsymbol{J} = 0 \arrow[r, ">> \delta = \star d \star"] & \delta \boldsymbol{J} = 0 \\
  & & \delta^2 = 0
\end{tikzcd}

\subsection*{BRST 量子化}
\[ Q^2 = 0, \quad Q\ket{\Psi} = 0, \quad \ket{\Psi} \sim \ket{\Psi} + Q\ket{\text{something}} \]

\section{自然偏爱简洁性}
物理公式如果太复杂，可能是：
\begin{enumerate}
  \item 它是个次级公式；
  \item 它是人为拼凑的；
  \item 它赖以成立的条件中存在错误或矛盾。
\end{enumerate}

越基本、普适的物理公式，形式越简洁。

\section*{Einstein 的名言}
\emph{“我想知道上帝是如何创造这个世界的。我对这个或那个现象，这个或那个元素的能谱不感兴趣。我要知道的是他的思想；其余的都是细节。”}

\end{document}
